\documentclass[a4paper]{amsart}

\usepackage{color}
\usepackage[colorlinks=true, citecolor=red]{hyperref}  %needs to come before amsrefs package to work with \cite{}*{}
\usepackage{amssymb,latexsym}
\usepackage[alphabetic]{amsrefs}
\usepackage[mathscr]{eucal}
\usepackage{pdfsync}
\usepackage{graphicx}
\usepackage{epstopdf}
\DeclareGraphicsRule{.tif}{png}{.png}{`convert #1 `dirname #1`/`basename #1 .tif`.png}
\usepackage[all]{xy}
\usepackage{titletoc}
%\usepackage{txfonts}   %  \fint  for integral average

%\usepackage{showkeys}  %options notref  notcite 

%\input{rgb.tex}

%\textcolor{DarkRed}{text in DarkRed}


\CompileMatrices

\theoremstyle{plain}
\newtheorem{theorem}{Theorem}[section]
\newtheorem{lemma}[theorem]{Lemma}
\newtheorem{proposition}[theorem]{Proposition}
\newtheorem{corollary}[theorem]{Corollary}

\theoremstyle{definition}
\newtheorem{definition}[theorem]{Definition}%[section]


\theoremstyle{remark}
\newtheorem{defn}[theorem]{Definition} 
\newtheorem{example}[theorem]{Example}
\newtheorem{remark}[theorem]{Remark}
\newtheorem{discussion}[theorem]{Discussion}

%\numberwithin{equation}{section}

%\newcommand{\abs}[1]{\lvert#1\rvert}

\DeclareMathOperator{\Lip}{Lip}
\DeclareMathOperator{\dist}{dist}
\DeclareMathOperator{\expected}{\mathbb{E}}
\DeclareMathOperator{\prb}{Prob}
\DeclareMathOperator{\spt}{spt}
\DeclareMathOperator{\lap}{\Delta}
\DeclareMathOperator{\Tr}{Tr}
\DeclareMathOperator{\co}{co}
\DeclareMathOperator{\diam}{diam}

\raggedbottom

%\addtolength{\textwidth}{2cm}
%\addtolength{\oddsidemargin}{-1cm}
%\addtolength{\evensidemargin}{-1cm} 
%\addtolength{\topmargin}{-1cm} 
%\addtolength{\textheight}{2cm}

 
%%% BEGIN DOCUMENT
\begin{document}

\title{Convexity \& Multifunctions}  
\author{John E.\ Hutchinson}
\date{\today}

 

\maketitle

\tableofcontents



\begin{abstract}  Motivation and ``geometric'' proofs of the material.  
\end{abstract}
%%%%%%%%%%%%%%%%%%%%%%%%%%%%%%%%
%%%%%%%%%%%%%%%%%%%%%%%%%%%%%%%%
%%%%%%%%%%%%%%%%%%%%%%%%%%%%%%%%
%%%%%%%%%%%%%%%%%%%%%%%%%%%%%%%%
\section{\textbf{A Result on Cones}}  
\noindent \emph{The following theorem for cones easily gives Theorem~\ref{mth}, which in turn easily implies both the Caratheodory and Shapley-Folkman Theorem.   The proof here is essentially that of \cite{Zho}, but giving the more general result in 
\cite{And} and the additional result in \cite{Sta}.}   




\bigskip The following Theorem is  geometrically clear if $d=2,3 $. 

For $  S \subset \mathbb{R}^d  $ the \emph{cone} $ c(S) $ is the set of all non-negative linear combinations
 \[
  c(S) = \sum_{i=1}^k \lambda_i x_i ,\quad  k \in \mathbb{N},\  x_i \in S, \  \lambda_i \geq 0   .
  \] 
The following says that $ c(S) $  is the union of the $ d $-dimensional sub-cones generated by  the subsets of $ S $ with cardinality $ d $.
 
\begin{theorem}\label{prco}
Suppose   $  S \subset \mathbb{R}^d  $ and $ x \in c(S) $.  Then 
\[
 x = \sum_{i=1}^k \lambda_i x_i ,\quad  \text{where } k\leq d,\  x_i \in S, \  \lambda_i \geq 0   .
 \]
\end{theorem} 


\begin{proof} We know that $ x $ can be expressed in the given form  for some $ k\in \mathbb{N}  $.  If $ k \leq  d $ we are done. 

Suppose  $ k > d $.  (We show we can reduce $ k $ by $ 1 $, the result then follows.)
W.l.o.g.\ take $ \lambda_i > 0 $ for all $ i $.
By linear dependence  there exist $ \mu_i $ not all zero such that 
$ \sum_{i=1}^k \beta _i x_i = 0 $. In particular,
\[
x =  \sum_{i=1}^{ k } (\lambda_i- \alpha \beta _i)  x_i
\]
for all $ \alpha \in \mathbb{R} $.  Choose $ \alpha >0 $ sufficiently small that $ \lambda_i- \alpha \beta _i \geq 0 $ for all $ i $, and 
 $ \lambda_j- \alpha \beta _j = 0 $ for some $ j $.

This reduces $ k $ by $ 1 $ and so gives the result.
\end{proof} 



%%%%%%%%%%%%%%%%%%%%%%%%%%%%%%%%
%%%%%%%%%%%%%%%%%%%%%%%%%%%%%%%%
%%%%%%%%%%%%%%%%%%%%%%%%%%%%%%%%
%%%%%%%%%%%%%%%%%%%%%%%%%%%%%%%%

\section{\textbf{Caratheodory's Theorem}}  

\noindent \emph{This section can be omitted as the result is also (easily) derived in Section~\ref{secSF}.}

\bigskip If $ P \subset \mathbb{R}^d  $ the \emph{convex hull} $ \co(P) $ of $ P $ is defined to be
\begin{itemize}
\item the intersection of all convex  sets containing $ P $;
\item equivalently, the set of all linear combinations
\[
\sum_{i=1}^k \lambda_i x_i ,\quad  k \in \mathbb{N},\  x_i \in P, \  \lambda_i \geq 0,\ \sum \lambda_i = 1   .
\]
\end{itemize}  
Caratheodory's theorem shows one can take $ k = d + 1 $.  That is, \emph{each $ x \in P $ is in a $ d    $-simplex with vertices  from $ P $}.  This is not surprising in dimensions 2 and 3.

\medskip Caratheodory's theorem follows easily from Theorem~\ref{prco}.  

\medskip
To see this, use the map
\begin{equation} \label{} 
x := (x^1, \dots, x^d) \in \mathbb{R}^d  \leftrightarrow  x^* := (x^1, \dots, x^d, 1) =: x^* \in \mathbb{R}^{d +1}.
\end{equation} 
For  $ P \subset \mathbb{R}^d  $   let $ P^* = \{x^* : x \in P \} $.  It  is  geometrically  clear that
\begin{equation} \label{} 
x \in \co (P) \text{ iff } x^* \in c (P^*).
\end{equation} 
The proof is immediate using the fact that the sum of the coefficients for points in the convex hull equals $ 1 $.  


 

\begin{theorem}[Caratheodory] Suppose $ P\subset \mathbb{R}^d  $ and $  x \in \co (P) $. Then
\begin{equation} \label{eq1} 
 x = \sum_{i=1}^{ k } \lambda_i x_i\  \text{ where } x_1,\dots , x_k \in P,\  k \leq d+ 1,\  \lambda_i > 0 ,\ \sum_{i=1}^{ k } \lambda_i = 1.
\end{equation} 
\end{theorem}  


\begin{proof}
If $ x \in \co (P) $ then $ x^* \in P^* $, and so by Proposition~\ref{prco}
\[
 x^* = \sum_{i=1}^k \lambda_i x^*_i ,\quad  \text{for some } k\leq d+1,\  x^*_i \in P^*, \  \lambda_i \geq 0   .
 \]
Moreover, $ \sum \lambda _i = 1 $ since the last component of both sides are equal, and from the left side must equal $ 1 $.
Taking the first $ d $ components gives the result.
\end{proof} 


The direct proof is also easy, and essentially just combines the proof of Theorem~\ref{prco} with the above proof.%
\footnote{\emph{Proof.} Suppose 
\[
x = \sum_{i=1}^k \lambda_i x_i ,\quad  k \geq d+2,\  x_i \in P, \  \lambda_i> 0,\ \sum \lambda_i = 1   .
\]
We will show we can reduce $ k $ by 1.  This gives the result.

By translation, without changing the $ \lambda_i $ we can  assume that $ x_1 = 0 $.  

Since the cardinality of $ \{x_2, \dots, x_k \} $ is at least $ d+1 $ it follows this set is linearly dependent. Hence for some $ \mu_2, \dots \mu_k $ not all zero and any $ \alpha \in \mathbb{R} $,  
\[
0 = \sum_{i=2}^k \mu_i x_i = \sum_{i=1}^k \mu_i x_i =  \sum_{i=1}^k \alpha \mu_i x_i,
\]
where $ \mu_1 $ is arbitrary.  Choose $ \mu_1 $ so $ \sum_{i=1}^k \mu_ i = 0 $.

It follows that 
\[
x =  \sum_{i=1}^k (\lambda_i - \alpha \mu_i ) x_i, \quad  \sum_{i=1}^k (\lambda_i - \alpha \mu_i ) = 1.
\]
Since each $ \lambda_i > 0 $, by increasing $ \alpha $ from $ 0 $ there is a value of $ \alpha > 0 $ such that 
$ \lambda_i - \alpha \mu_i \geq 0 $ for all $1\leq i \leq k $, and $ \lambda_i - \alpha \mu_j = 0 $ for some $1\leq j \leq k $.

This gives   \eqref{eq1} with $ k $ reduced by 1 and so completes the proof.
}


%%%%%%%%%%%%%%%%%%%%%%%%%%%%%%%%
%%%%%%%%%%%%%%%%%%%%%%%%%%%%%%%%
%%%%%%%%%%%%%%%%%%%%%%%%%%%%%%%%
%%%%%%%%%%%%%%%%%%%%%%%%%%%%%%%%

\section{\textbf{Shapley-Folkman Theorem}} \label{secSF}  
\noindent \emph{The proof here is essentially that of \cite{Zho}, but gives the more general result in 
\cite{And}.  Caratheodory and Shapley-Folkman follow immediately.}

\bigskip

\begin{theorem}\label{mth}
Suppose $ x \in co (A_1 + \dots + A_N) $ where $ A_i \subset \mathbb{R}^d  $.  Then
\begin{equation} \label{} 
x = \sum_{i=1}^N \sum_{j=1} ^{m_i} \lambda _{ij} a_{ij} ,
\quad \text{where }  a_{ij}\in A_i,\  \sum_{i=1}^N m_i \leq N + d, \ \lambda _{ij} > 0 ,\ \sum_{j=1}^{m_i} \lambda _{ij} = 1.
\end{equation} 
\end{theorem} 

\begin{remark}\label{rem1}
It is immediate that 
\[
x = \sum_{j=1}^m \lambda _j \sum_{i=1}^N a_{ij} = \sum_{i=1}^N \sum_{j=1} ^{m } \lambda _{j} a_{ij} ,
\quad \text{where }  a_{ij}\in A_i,  \ \lambda _{j} > 0 ,\ \sum_{j=1}^{m } \lambda _{j} = 1.
\]
So the point is that \emph{one can bound the number of non-zero terms by $ N+ d$} (at the cost of replacing $ \lambda _j $ by $ \lambda _{ij} $).
\hfill $\square$  \end{remark} 

 
 \begin{corollary}[Caratheodory Theorem]
 Suppose $ P\subset \mathbb{R}^d  $ and $  x \in \co (P) $. Then
\begin{equation} 
 x = \sum_{i=1}^{ k } \lambda_i x_i\  \text{ where } x_1,\dots , x_k \in P,\  k \leq d+ 1,\  \lambda_i > 0 ,\ \sum_{i=1}^{ k } \lambda_i = 1.
\end{equation} 
\end{corollary}  

\begin{proof}
Let $ N=1 $.
\end{proof}

For the next corollary think of the case $ N \gg d $.
\begin{corollary}[Shapley-Folkman Theorem]  Suppose  $ x \in \co (A_1 + \dots + A_N) $ where $ A_i \subset \mathbb{R}^d  $.  

Then
\begin{equation} \label{sfd}
x = a_1 + \dots + a_N \in \sum_{i\notin J} A_i + \sum _{i\in J} \co(A_i) =  \sum_{i\notin J} A_i + \co \left(\sum _{i\in J}  A_i \right)
\end{equation}
where  $ J \subset N $ depends on $ x $,    $ |J|\leq d $, $ a_i \in A_i $ if $ i \notin J $ and $ a_i \in \co (A_i ) $ if $ i \in J $.

Trivially,  $ A_1 + \dots + A_N \subset \co (A_1 + \dots + A_N)$. 
If $ N \geq d $ then there exists $ y \in A_1 + \dots + A_N $ such that 
\begin{equation} \label{}
 |x-y| \leq \sqrt{d} \,L \text{ where } L := \max_{1 \leq i \leq N} \diam A_i .
 \end{equation} 
\end{corollary} 

\begin{proof} In theorem~\ref{mth}, since $ m_i \geq 1 $ and   $ \sum_{i=1}^N m_i \leq N + d $, $ m_i = 1 $ for at least $ N- d $ of the $ i $.  

The last equation in~\eqref{sfd} follows from~\eqref{scs}. 

\medskip For the last part choose $ y = b_1 + \dots +   b_N \in A_1 + \dots + A_N $, where 
\[
b_i = \begin{cases} a_i   & i \notin J \\  \text{any } a_i^*\in A_i  & i \in J \end{cases}.
\]
Then
\[
|x-y| \leq \sum_{i\in J} |a_i^* - a_i| \leq \sqrt{d}\, L .
\]
\end{proof}

\begin{remark} The significance of the last part of the Shapley-Folkman theorem is that it shows the sum of a large number $ N $ of (not necessarily convex) sets of bounded diameter becomes more convex as $ N $ increases.  The distance between $ A_1 + \dots + A_N  $  and $ \co (A_1 + \dots + A_N) $   is bounded independent of $ N $ but the diameter of the sum is generically of   order $ N $.  After normalising by the factor $ N^{-1} $ the difference between the sum and its convex hull is of order $ O(1/N) $.

An example is $ A_i = \{ (0,0), (0,1), (1,0), (1,1)\} \subset \mathbb{R}^2 $ for all $ i $.  Then $ A_1 + \dots + A_N  $ is the set of integer grid points in $ [0,1]\times [0,1] $ and $ \co (A_1 + \dots + A_N)  =  [0,N]\times [0,N] $.
\hfill $\square$  \end{remark} 




We now prove Theorem~\ref{mth}.
\begin{proof}[Proof of Shapley-Folkman Theorem]  As in Remark~\ref{rem1} let 
\[
 x =   \sum_{j=1}^m \lambda _j \sum_{i=1}^N a_{ij} = \sum_{i=1}^N \sum_{j=1} ^{m } \lambda _{j} a_{ij} ,
\quad \text{where }  a_{ij}\in A_i,  \ \lambda _{j} > 0 ,\ \sum_{j=1}^{m } \lambda _{j} = 1.
\]
Define the following points in $ \mathbb{R}^{d+N}  $ for $ j = 1,\dots, N $.
\begin{align*}
x^* &= (x,1,\dots, 1) \\
a_{1j}^* &= (a_{ij},1,0,\dots, 0 ) \\
            &\ \vdots \\
a_{Nj}^* &= (a_{ij},0,0,\dots, 1 )  
\end{align*}  
Then using  $  \sum_{j=1}^{m } \lambda _{j} = 1 $, 
\[
 x^*   = \sum_{i=1}^N \sum_{j=1} ^{m } \lambda _{j} a^*_{ij}  
%\quad \text{where }  a_{ij}\in A_i,  \ \lambda _{j} > 0 ,\ \sum_{j=1}^{m } \lambda _{j} = 1.
\]

By Theorem~\ref{prco} applied in  $ \mathbb{R}^{d+N}  $ with $ S = \{ a^*_{ij}: 1\leq i \leq N, 1 \leq j \leq m    \} $, after relabelling the $ a_{ij} \in A_i $,
\[
 x^*  = \sum_{i=1}^N \sum_{j=1} ^{m_i } \lambda _{ij} a^*_{ij} ,
\quad \text{where }  \sum_{i=1}^N m_i \leq d + N,\  a_{ij}\in A_i,  \ \lambda _{ij} > 0 .
\]
Equating each of the last $ N $ components on either side gives, for $ 1 \leq i \leq N $,
\[ 1 = \sum_{j=1}^{m_i} \lambda _{ij} . \]
Taking the first $ d $ components now gives the required result.
\end{proof} 

\begin{remark}\emph{Supporting hyperplanes.}
For any $ S \subset \mathbb{R}^d  $, $ \co(S) $ is the intersection of all \emph{open} halfspaces containing (supporting) $ S $. If $ S $ is closed then one can use closed half spaces containing $ S $, and in any case using closed hyperspaces gives $ \overline{\co} (S) $, the closure of $ \co (S) $. 

For any   unit vector $ p \in \mathbb{R}^d  $ and $ \alpha \in \mathbb{R} $ the associated (closed) half space and oriented hyperplane are
\[
  \{ x : x \cdot p \leq \alpha \},  \quad 
  \{ x : x \cdot p = \alpha \} \text{ with orientation $p $ (not $ -p $)},
  \]
and denoted by $ \alpha p $ in either case.   

If $ S \subset \mathbb{R}^d  $ and $ S $ is a subset of the halfspace $ \alpha p $ then we say $ S $ is \emph{supported} by the corresponding halfspace or hyperplane.      

 If $ S $ is \emph{closed and bounded} then in each   unit direction $ p $ there is a \emph{least} $ \alpha $ such that $ S $ is supported by $ \alpha p $. 
We call this the   \emph{minimal supporting halfspace} or \emph{minimal supporting hyperplane} for $ S $. Note that in this case  \emph{$ \co(S )$ is the intersection of the minimal supporting halfspaces for $  S $}.
\hfill $\square$  \end{remark} 

\begin{theorem} If $ S_i \subset \mathbb{R}^d  $ then 
\begin{equation} \label{scs} 
\co(S_1 + \dots + S_N) = \co (S_1) + \dots + \co (S_N) .
\end{equation} 
\end{theorem}

\begin{proof}
The inclusion ``$ \subset $'' is immediate from Remark~\ref{rem1}.

\medskip 
For ``$ \supset $'' first note that a sum of convex sets is convex, and so the right side of \eqref{scs} is  convex.
 For simplicity assume the $ S_i $ are closed and bounded (the case we need in   applications).  Then both sides are closed, bounded and convex.
 
Fix a unit vector $ p \in \mathbb{R}^d  $. Suppose $ \alpha _i p $ is the minimal supporting hyperplane for $ S_i $, and hence for $ \co (S_i )$. 

Then $ (\sum_i \alpha _i ) p $ is the minimal supporting hyperplane for $ S_1 + \dots +S_N $ and $ \co (S_1)+ \dots + \co(S_N) $ since
\begin{align*} 
x_i \cdot p \leq \alpha _i &\Longrightarrow \left(\sum x_i\right)  \cdot p \leq  \sum\alpha _i  ,\\
x_i \cdot p = \alpha _i &\Longrightarrow \left(\sum x_i\right)  \cdot p = \sum \alpha _i  .
\end{align*}  
Thus both  $ S_1 + \dots +S_N $ and $ \co (S_1)+ \dots + \co(S_N) $ have the same closed convex hull, which gives \eqref{scs}.
\end{proof} 

 


%%%%%%%%%%%%%%%%%%%%%%%%%%%%%%%%
%%%%%%%%%%%%%%%%%%%%%%%%%%%%%%%%
%%%%%%%%%%%%%%%%%%%%%%%%%%%%%%%%
%%%%%%%%%%%%%%%%%%%%%%%%%%%%%%%%

\section{\textbf{Lyapunov Theorem}}  
\noindent \emph{The proof follows \cite{Tar} and is a relatively elementary consequence of the Shapley-Folkman theorem.}

\begin{remark} \emph{Measure theoretic preliminaries.}   \label{mtp}
\begin{enumerate}
\item Fix a finite positive measure space $ (X, \Sigma, \mu) $. 

 \item An \emph{atom} is a set $ A $ such that 
 \[ \mu(A) > 0\quad  \text{and} \quad  B\subset A \Longrightarrow \mu(B) = 0 \text{ or } \mu(B) =\mu(A).
 \]
  (Whenever we write $ \mu(E) $ it is assumed $ E \in \Sigma $.) We can regard atoms as points by taking the isomorphic measure space given by the equivalence relation under which two points are equivalent iff they are equal or belong to the same atom.

 
\item $ \mu $ is defined to be \emph{atomless} if there are no atoms. That is,
\[
 \mu(A) > 0 \Longrightarrow \exists B\subset A \text{ s.t.\ } 0 < \mu(B) < \mu(A).
 \]
 
 \item If $ \mu $ is atomless then
 \[
 0 < \alpha < \mu(A) \Longrightarrow \exists B \text{ s.t.\ } \mu(B) = \alpha .
 \]
\begin{proof} (See \cite{Hal}.)  By beginning with $ A $ and repeatedly applying the definition of atomless,   taking either the set constructed or its complement in the set previously constructed, we obtain a subset of $ A $ of arbitrarily small positive measure.

Thus   $\exists B_1\subset  A $ with $ 0 < \mu(B_1) \leq \alpha $.   If $   \mu(B_1) = \alpha $ we are done.  If not, $\exists  B_2 \subset A \setminus B_1 $ with $0<  \mu(B_2) \leq  \alpha - \mu(B_1) $. In this way we construct a well ordering of disjoint sets $ B_1, B_2 , \dots $  which are subsets of $ A $ and have \emph{positive} measure.  This w.o.\ must have countable length, and the union gives the required set $ B $.
 \end{proof}  
 
 \item $ \mu = (\mu_1, \dots, \mu_d) $ is a \emph{vector measure} on $ (X,\Sigma) $ with values in $ \mathbb{R}^d  $ if each $ \mu_i $ is a finite signed measure on $ (X,\Sigma) $.
 
 $ \mu $ is said to be \emph{atomless} if each variation measure $ |\mu_i | $ is atomless.
 
 \item The \emph{total variation measure} is $ |\mu| = |\mu_1| + \dots + | \mu_d| $.
  By Radon Nikodym there is a $ |\mu| $ integrable function $ f : X \to \mathbb{R}^d  $ such that 
 \begin{equation} \label{}
\frac{d\mu }{d |\mu| } = f , \quad  \text{hence } \mu(A) = \int_A f \, d|\mu| \ \ \forall A \in \Sigma.
 \end{equation}
 \end{enumerate} 
  \hfill $\square$  \end{remark} 


 
 
 \begin{remark} \label{rvvm}  \emph{The range of vector-valued measures.}
 
 \begin{enumerate}
 \item The \emph{range} 
 \[
  \mathcal{R} (\mu) = \mathcal{R} (X) := \{ \mu(E) : E \subset X \} \subset \mathbb{R}^d  
  \]
   of $ \mu $ is its range as a map $ \mu : \Sigma \to \mathbb{R}^d $.  That is, the set of values of $ \mu $ over all possible ``coalitions'' (in economics).
  For fixed $ \mu $ and $ A \in \Sigma$, the \emph{range of $ \mu $ restricted to $ A $} is 
 \[
 \mathcal{R} (A) :=  \{ \mu(E) : E \subset A \} .
 \]

   It is clear that 
 if $ \{A_1,\dots, A_N\} $ is a (measurable) partition of $ X $ then 
 \[
\mathcal{R} (X) = \mathcal{R} (A_1) + \dots + \mathcal{R} (A_N).
 \]
 
 
 \item If $ \mu $ is positive real and atomless then, by Remark~\ref{mtp}\,(4), $ \mathcal{R} (\mu) = [0,\mu(X)] $.
 
 If $ \mu $ is real and atomless it follows that $ \mathcal{R} = [-\mu^-(X), \mu^+(X)] $, where $ \mu = \mu^+ - \mu^- $ is the standard Hahn-Jordan decomposition.
 
For such real and atomless $ \mu $, if $ \nu = \mu + \alpha _1  \delta _{x_1} + \dots +  \alpha _n \delta _{x_n} $ then it follows that
 \[
 \mathcal{R} (\nu) =  [-\mu^-(X), \mu^+(X)]  + \{0,\alpha _1\} + \dots + \{0,\alpha_n\}.
 \]
 \end{enumerate} 
 \hfill $\square$  \end{remark}
 
 \begin{remark}
Here is   motivation for the following theorem.

Let $ X_1,\dots, X_N $ be a partition of $ X $ with $ N $ large, 
\[
 |\mu| (X_i)= N^{-1}|\mu|(X)  , \quad  \mu(X_i) \approx |\mu|(X_i)\, \widehat{\mu}_i \text{ where $ \widehat{\mu}_i $ is a unit vector} .
 \]
 Then
 \begin{align*}
 \mathcal{R} (\mu) &\approx \bigg\{ \mu \bigg( \bigcup_{i\in J} X_i \bigg) : J \subset \{1,\dots,N\} \bigg\}
    \quad \text{in the Hausdorff metric} \\
    &= \bigg\{ \sum_{i\in J} \mu(X_i) : J \subset \{1,\dots,  N \} \bigg\} \\
    &= \sum_{i=1}^N \{ 0, \mu(X_i)\} \quad \text{NOTE} \\
    &\approx \sum_{i=1}^N [ 0, \mu(X_i)] \quad \text{by Shapley-Folkman} \\
    &\approx  |\mu|(X)\, N^{-1} \sum_{i=1}^N [ 0, \widehat{\mu}_i ] .
 \end{align*} 
 
  This also indicates that $ \mathcal{R} (\mu) $ is the limit of zonotopes, which is   a theorem.
\hfill $\square$  \end{remark} 
  
 
 The proof of the following is essentially from \cite{Tar}.  However, I proceed directly to the proof of the Lyapunov theorem without going via the result on multifunctions in Theorem~\ref{mfs}.  The proof of closedness is simpler --- the proof in~\cite{Tar} essentially reproves the Hahn-Jordan decomposition via the Radon Nikodym theorem.
 
 \begin{theorem} \label{liap} If the vector-valued measure $ \mu $ is atomless then $ \mathbb{R} (\mu) $ is compact and convex.
 \end{theorem}
 
 
 \begin{proof}  Consider  $ (X,\Sigma, \mu) $ where $ \mu $ is an $ \mathbb{R}^d   $ valued measure.  When we write $ E\subset X $ we also assume $ E $ is measurable, that is  $ E \in \Sigma$. 
 
 \bigskip
 \noindent \emph{Boundedness}.  If $ E \subset X $ then
 \begin{align*}
 |\mu(E)| &=  \Big| \big( \mu_1(E),\dots, \mu_d (E)\big) \Big| \\
    &=  \Big| \big( |\mu_1 (E)|,\dots, |\mu_d (E)|\big) \Big| \\
   &\leq \Big| \big(| \mu_1|(E),\dots, |\mu_d| (E)\big) \Big|  \\
   &=: K < \infty.
   \end{align*} 
   
\bigskip
\noindent \emph{Convexity.}  Suppose $ a\in \co \mathcal{R} (\mu) $.  We W.T.S.\ $ a \in \mathcal{R} (\mu) $, i.e.\ $ a = \mu (E) $ for some $ E\subset  X $.

Since $ |\mu | $ is atomless, $ \exists $ a (measurable) partition $ X_1,\dots, X_{2d} $ of $ X $ such that 
$ |\mu|(X_i) = 2^{-2d}\, |\mu| (X ) $.

Using the notation of Remark~\ref{rvvm}\,(1), let $ A = \mathcal{R} (X) \subset \mathbb{R}^d  $ and $ A_i = \mathcal{R} (X_i) \subset \mathbb{R}^d  $.  Then, using additivity of measures,
\begin{align*}
A &= A_1 + \dots + A_{2d} \\
a \in \co (A) &= \co( A_1 + \dots + A_{2d}).
\end{align*} 

By the Shapley-Folkman theorem,
\[
a = a_1 + r_1, 
\]
where, for some $ J \subset \{1,\dots, 2d\} $ with $ |J|= d$, 
\begin{gather*}
a_1 \in \bigcup_{i\notin J} A_i = \mathcal{R} \bigg( \bigcup_{i \notin J} X_i \bigg) = \mathcal{R} (S_1)
,\quad  S_1 := \bigcup_{i\notin J}X_i,\    |\mu|(S_1) = 2^{-1} |\mu|(X),            \\
r_1 \in \co \bigg(\bigcup_{i\in J} A_i \bigg)= \co \mathcal{R} \bigg( \bigcup_{i \notin J} X_i \bigg) = \co \mathcal{R} (X\setminus S_1),
\quad  |\mu|(X\setminus S_1) = 2^{-1} |\mu|(X).
\end{gather*}
 
Iteratively applying this argument with $ X $ replaced by $ X\setminus S_1 $, $ X \setminus (S_1\cup S_2) $, \dots , we obtain
\begin{equation} \label{air} 
a = a_1 + \dots + a_n + r_n,
\end{equation} 
where 
\begin{gather*}
a_i \in \mathcal{R} (S_i), \quad S_1,\dots, S_n \text{ are disjoint}, \\
|\mu|(S_i) = 2^{-i}\, |\mu|(X), \quad |\mu|\bigg(X\setminus \bigcup_{i=1}^n S_i\bigg) = 2^{-n} |\mu|(X).
\end{gather*}
It follows that
\begin{gather*}
a_i  = \mu(E_i ) \text{ for some } E_i  \subset S_i, \quad \therefore |a_i| \leq |\mu|(S_i) \leq 2^{-i} |\mu|(X),\\
r_n \in \co \bigg\{ \mu(E ) : E  \subset X \setminus \bigcup_{i=1}^nS_i\bigg\}, \quad 
             \therefore |r_n| \leq |\mu|\bigg(X\setminus \bigcup_{i=1}^n S_i\bigg) \leq 2^{-n} |\mu|(X).
\end{gather*}

Hence from \eqref{air},
\[
a  = \sum_{i\geq 1} a_i  
 = \sum_{i\geq 1} \mu(E_i) = \mu \bigg( \bigcup_{i\geq 1} E_i \bigg) \in \mathcal{R} (X).
 \]
 
 \bigskip 
 \noindent \emph{Closedness.}  
 Suppose $ a \in \overline{\mathcal{R} (\mu)} \setminus \mathcal{R} (\mu) $.  
 
 By convexity there is a hyperplane containing $ a $ such that one of the corresponding closed hyperspaces contains $ \mathcal{R} (\mu) $.  By rotating coordinates in $ \mathbb{R}^d  $ we can assume
 \begin{align*}
 a =(a^1, \dots, a^d) &= \lim_{i\to \infty} \mu (X_i ),  \quad X_i \subset X,\\
 a^d &= \sup \{ \mu_d(Y) : Y \subset X \}.
 \end{align*} 
 
 Using Radon Nikodym let 
 \[
 \mu = f \, |\mu|, \quad f = (f^1,\dots, f^d).
 \]
 Let
 \[
 P = \{ x\in X : f^d(x) > 0 \}, \   Z = \{ x\in X : f^d(x) = 0 \}, \  N = \{ x\in X : f^d(x) < 0 \}.
 \]
 
Then
 \[
 \mu_d(P) = a^d,\quad   \mu^d(Z) = 0,\quad \mu^d(N) \leq 0, \quad  \mu^d(N)  < 0 \text{ if } |\mu|(N)>0 .
 \]
It follows  
\[
\mu^d (X_i \cap N )  \to 0, \quad 
\therefore |\mu| (X_i \cap N )  \to 0, \quad 
\therefore  \mu  (X_i \cap N )  \to 0.
\]
 Similarly
 \[
\mu^d (X_i \cap P )  \to \mu^d(P), \quad 
\therefore |\mu| (X_i \cap P )  \to |\mu|(P)|\text{ and }
   \mu   (X_i \cap P )  \to \mu(P).
\]

 
 Hence
 \[
 a = \lim \mu(X_i) = \lim \mu(X_i\cap P) + \lim \mu(X_i\cap Z) = \mu(P) + \lim \mu(X_i \cap Z ).
 \]
 Since the range of $ \mu $ restricted to $ Z $ has one less dimension, we can use induction to deduce the result.
 \end{proof}  
 


 

%%%%%%%%%%%%%%%%%%%%%%%%%%%%%%%%
%%%%%%%%%%%%%%%%%%%%%%%%%%%%%%%%
%%%%%%%%%%%%%%%%%%%%%%%%%%%%%%%%
%%%%%%%%%%%%%%%%%%%%%%%%%%%%%%%%
\section{\textbf{Multifunctions and Measure}}

 
\begin{remark} \emph{Preliminaries}
\begin{enumerate}
\item A function $ F $ is a \emph{multifunction} from $ X $ to $ \mathbb{R}^d  $ if $ F(x) \subset \mathbb{R}^d  $ for all $ x \in X $.
We write $  F:X \rightrightarrows \mathbb{R}^d $.

\item $ f:X \to \mathbb{R} $ is an \emph{integrable selection} for  $ F $ if $ f $ is integrable and $ f(x) \in F(x) $ for all $ x\in X $.
 
\item $ \int_X F := \{ \int f : f \text{ is an integrable selection for } F \}$. 
\end{enumerate}
\hfill $\square$  \end{remark}

 
\begin{remark} \label{nfo}  \emph{New integrable selections from old.}  Let $ \mathcal{L} _F $ denote the class of integrable selections from $ F $.  

If $ f_1, f_2 \in \mathcal{L} _F $ then $ \mathcal{X} _E f_1 + \mathcal{X} _{E^c} f_2 \in  \mathcal{L} _F $ for any $ E \in \Sigma $.
By taking $ F(x) = \{ f_1(x), f_2(x) \} $ we see this gives \emph{all} integrable selections in $ \mathcal{L} _F $ which follow from just the knowledge that $ f_1, f_2 \in \mathcal{L} _F $.
\hfill $\square$  \end{remark}  
 
The proof of Lyapunov's theorem easily generalises to give the following:
 
\begin{theorem}\label{mfs}
  If the finite positive  measure $ \mu $ is atomless and $  F :X\rightrightarrows \mathbb{R}^d $ then $ \int F\, d\mu  $ is convex.
 \end{theorem} 
 
 \begin{proof}[Direct Proof Outline]  The proof is essentially as for convexity in Theorem~\ref{liap}.
 
 Suppose $ f,g \in F$.  Define the multifunction $ G $ by
 \[
    G(x) = \{ f(x), g(x) \} \subset F(x).
    \]
    It is sufficient to show $ \int G \, d\mu $ is convex.
    
 Similarly to before, take a partition $ X_1,\dots, X_{2d} $ of $ X $ but now with $ \mu(X_i) = 2^{-2d} \mu(X) $. As before, suppose
 \[
 a \in A:= \int G\, d\mu = \sum_{i=1}^{2d}\int_{X_i}G\, d\mu =: \sum_{i=1}^{2d} A_i,
 \]
 where $ A, A_i \subset \mathbb{R}^d  $.
 
Again as before, by appropriately iterating the argument, we obtain
\[
a = \sum_{i=1}^\infty a_i \in \sum_{i=1}^\infty \int_{S_i} G = \int_{\bigcup_{i=1}^\infty S_i} G = \int _X G,
\]
with equality since $ \mu\big(   \bigcup_{i=1}^\infty S_i\big) = \mu(X) $ by the construction.
 \end{proof}

 \begin{proof}[Proof from Lyapunov Theorem]  Suppose $ f_1,f_2 \in \mathcal{L}_F $ and $ 0 < \lambda < 1 $.  
 
 To show $ \int F  $ is convex it is sufficient (and necessary) to show 
 \[
 \int f   = \lambda \int f_1 + (1- \lambda ) \int f_2 .
 \]
 for some $ f \in \mathcal{L} _F $.
 
 From Remark~\ref{nfo} it is sufficient (and necessary) to find such an $ f $ of the form $ 
\mathcal{X} _E f_1 + \mathcal{X} _{E^c} f_2 $ for some $ E \in \Sigma $. 
That is, we require
 \begin{align*}
  \lambda \int f_1 + (1- \lambda ) \int f_2  &=  \int \big( \mathcal{X} _E f_1 + \mathcal{X} _{E^c} f_2  \big) \\
    \text{i.e.\ } &= \int_E f_1 + \int  f_2 - \int_{E^c} f_2 ,   \\
    \text{i.e.\ } \left( \int_E f_1, \int_E f_2\right) &= \lambda \left(\int f_1, \int f_2 \right) .
 \end{align*}
 
 The existence of such an $ E $ follows by applying the Lyapunov theorem to the atomless  $ \mathbb{R}^{2d} $-valued measure $ \nu $ defined by
 \[
 \nu(E) =  \left( \int_E f_1, \int_E f_2\right),
 \]
 and noting $ (0,0), \left(\int f_1, \int f_2 \right) \in \mathcal{R} (\nu) $.
  \end{proof}
 
 
 %%%%%%%%%%%%%%%%%%%%%%%%%%%%%%%%
%%%%%%%%%%%%%%%%%%%%%%%%%%%%%%%%
%%%%%%%%%%%%%%%%%%%%%%%%%%%%%%%%
%%%%%%%%%%%%%%%%%%%%%%%%%%%%%%%%
\section{\textbf{Multifunctions and Continuity}}  \emph{In the following, $ X $ and $ Y $ are metric spaces.}

\begin{definition} A multifunction $ F:X \rightrightarrows Y $  is \emph{upper semi continuous}
(or \emph{upper hemi continuous}) at $ x \in X $ if for every open $ V \supset F(X) $ there is an open $ U \ni x $ s.t.\ $ F(z) \subset V $ for every $ x \in U $.  We write ``u.s.c'' or ``u.h.c.''.
\end{definition}

It follows that a \emph{function} $ f :X \to Y $ is continous iff it is u.s.c.\ as a multifunction.

\begin{definition} A multifunction $ F:X \rightrightarrows Y $ is \emph{closed} at $ x \in X $ if 
\[
x_n \to x,\ y_n \in F(x_n),\ y_n \to y \Longrightarrow y \in F(x).
\]
\end{definition} 

An example of $ F:X \rightrightarrows Y $ which is compact valued and is closed at every $ x $, but is not u.s.c.,  is $F= f:[0,1]\to \mathbb{R} $ where $ f(0) = 0 $, and $ f(x) = x^{-1} $ if $ x\in(0,1] $.

However

\begin{proposition} 
If $ Y $ is compact and $   F:X \rightrightarrows Y $ is compact valued at $ x $, then $ F $ is u.s.c.\ iff it is closed at $ x $.
\end{proposition}

\begin{proof} Straightforward. \end{proof} 



%%%%%%%%%%%%%%%%%%%%%%%%%%%%%%%%
%%%%%%%%%%%%%%%%%%%%%%%%%%%%%%%%
%%%%%%%%%%%%%%%%%%%%%%%%%%%%%%%%
%%%%%%%%%%%%%%%%%%%%%%%%%%%%%%%%
\section{\textbf{Limits of Sets}}
  \emph{In the following, $ X $ is a  metric space unless noted otherwise.}

\begin{definition} Suppose $ X_n \subset X $ for $ n\geq 1 $. Then
\begin{align*}
x \in \limsup X_n &\text{ iff } \exists   n_k \to \infty\ \& \  \exists x_k \in X_{n_k} \text{ with }  x_{n_k}\to x ,  \\
x \in \liminf X_n &\text{ iff }   \exists x_n  \in X_{n } \text{ with }   x_n \to x.
\end{align*} 
\end{definition} 
 
For an arbitrary set $ X $ (not necessarily with  metric)  
 \begin{align*}
\limsup X_n &= \bigcap _{n\geq 1} \bigcup_{k\geq n} X_k = \{ x : x \in X_n \text{ i.e.}\} ,\\
\liminf X_n &= \bigcup _{n\geq 1} \bigcap_{k\geq n} X_k = \{ x : x \in X_n \text{ eventually}\} .\\
 \end{align*} 
The two definitions in each case agree if $ X $ is given the discrete metric.

\bigskip
It is easy to check that
\begin{align*}
\limsup X_n &= \limsup \overline{X_n} =\overline{\limsup X_n},\\
\liminf X_n &= \liminf \overline{X_n} =\overline{\liminf X_n}.\\
\end{align*} 







%%%%%%%%%%%%%%%%%%%%%%%%%%%%%%%%
%%%%%%%%%%%%%%%%%%%%%%%%%%%%%%%%
%%%%%%%%%%%%%%%%%%%%%%%%%%%%%%%%
%%%%%%%%%%%%%%%%%%%%%%%%%%%%%%%%
\section{\textbf{More}} 
 
For a discussion and quick proof 
\begin{enumerate}
\item
of Fatou's Lemma in $ n $-dimensions,
\item that integration preserves upper semicontinuity,
\end{enumerate}
see \cite{Aum} and \cite{RY}.
 

%%%%%%%%%%%%%%%%%%%%%%%%%%%%%%%%
%%%%%%%%%%%%%%%%%%%%%%%%%%%%%%%%
%%%%%%%%%%%%%%%%%%%%%%%%%%%%%%%%
%%%%%%%%%%%%%%%%%%%%%%%%%%%%%%%%

\begin{thebibliography}{9} 

\bib{And}{misc}{
  author={Anderson, Robert}
  title={Nonconvex Preferences and Approximate Equilibria, Econ 201B}
  date={2010}
  note={Lecture Notes online}
  }
  
  \bib{Aum}{article}{
   author={Aumann, Robert J.},
   title={An elementary proof that integration preserves
   uppersemicontinuity},
   journal={J. Math. Econom.},
   volume={3},
   date={1976},
%   number={1},
   pages={15--18},
%   issn={0304-4068},
%   review={\MR{0403571 (53 \#7382)}},
}
		
  \bib{Hal}{article}{
   author={Halmos, Paul R.},
   title={On the set of values of a finite measure},
   journal={Bull. Amer. Math. Soc.},
   volume={53},
   date={1947},
   pages={138--141},
%   issn={0002-9904},
%   review={\MR{0020124 (8,506d)}},
}

\bib{RY}{report}{
  author={Rustichini,Aldo},
  author={Yannelis, Nicholas C.},
  title={An Elementary Proof of Fatou's Lemma in Finite Dimensional Spaces},
  date={1986},
  number={237},
  organization={Center for Economic Research, University of Minnesota},
  type={Discussion paper},
  }
  
  

\bib{Sta}{misc}{
  author={Starr, Ross}
  title={Shapley-Folkman Theorem}
  date={2008}
  note={Lecture Notes online}
  }

\bib{Tar}{article}{
   author={Tardella, Fabio},
   title={A new proof of the Lyapunov convexity theorem},
   journal={SIAM J. Control Optim.},
   volume={28},
   date={1990},
%   number={2},
   pages={478--481},
%   issn={0363-0129},
%   review={\MR{1040471 (91c:28009)}},
%   doi={10.1137/0328026},
}

\bib{Zho}{article}{
   author={Zhou, Lin},
   title={A simple proof of the Shapley-Folkman theorem},
   journal={Econom. Theory},
   volume={3},
   date={1993},
   number={2},
   pages={371--372},
%   issn={0938-2259},
%   review={\MR{1216029}},
%   doi={10.1007/BF01212924},
}	



\end{thebibliography}

\end{document} 

